\documentclass[12pt,a4paper]{article}
\usepackage[utf8]{inputenc}
\usepackage[T1]{fontenc}
\usepackage{mathptmx}
\usepackage[margin=2cm]{geometry}
\usepackage{amsmath}
\usepackage{enumitem}
\usepackage{hyperref}
\hypersetup{colorlinks=false, pdfborder={0 0 0}}

\title{\textbf{Threshold-Aware Similarity Caching for\\
Named Data Networks:\\
Bounding Retry Delay under Congestion using ns-3}}
\author{Somyak Priyadarshi Mohanta (241CS257)\\
Pranav Shaji (241CSXXX)\\
Chris Tony (241CSXXX)\\
\textbf{Department of Computer Science and Engineering}\\
\textbf{National Institute of Technology, Karnataka}}
\date{\today}

\begin{document}
\maketitle

\section*{Abstract}

\subsection*{Background}
Content caching is one of the main techniques used to deliver content quickly and to reduce the
traffic carried by a network. Content Delivery Networks, mobile edge caches and the in-network
caches of Named Data Networking (NDN) all follow the same rule: a cache serves a request only when
it holds the exact object that was asked for [1]. Similarity caching relaxes this rule. On a miss,
the cache may return the most similar object it holds, which is useful for applications such as
image retrieval, recommendation systems and semantic caching of large language model queries
[2], [3]. The user accepts the returned object only if it is similar enough, and otherwise sends
the request again.

\subsection*{State-of-the-Art}
Nakamura and Kamiyama recently analysed similarity caching on a general cache network and
introduced the similar-content delivery delay, a user-level metric that includes the time spent on
repeated requests [4]. In their model each cache answers with a fixed probability $\alpha$, and the
network is assumed to have infinite bandwidth, no queuing and no packet loss. The results are
obtained only through numerical evaluation on a tandem topology of five caches. Other studies
consider a single similarity cache or optimise content placement in networks of similarity caches
[2], [5], but none of them evaluates the request and retry behaviour at the packet level.

\subsection*{Issues and Challenges}
Several challenges exist in making similarity caching work in a real multi-hop network:
\begin{itemize}
  \item \textbf{Correlated Retries}: The analysis treats every retry as an independent draw of the
  cache contents. In practice the contents do not change within the retry interval, so a user who
  rejects an answer keeps receiving the same one. Re-evaluating the model on the same topology with
  persistent cache contents, we found that the mean delay is underestimated by about two times at
  $\alpha = 0.5$ and by more than sixty times at $\alpha = 0.9$.
  \item \textbf{Blind Caches}: A cache does not know the similarity that the user will accept, so it
  can return an object that is certain to be rejected.
  \item \textbf{NDN Name Matching}: NDN delivers a Data packet only when its name matches the
  Interest, so returning a different but similar object needs changes to the packet format and
  the forwarder.
  \item \textbf{Congestion and Loss}: The benefit of answering from a nearby cache depends on the
  queuing delay and loss on the upstream path, which the existing model does not capture.
\end{itemize}

\subsection*{Objectives}
This project aims to design and implement threshold-aware similarity caching for NDN and to
evaluate it in ns-3 using the ndnSIM module [6]. The specific objectives are:
\begin{enumerate}
  \item Implement a similarity-based Content Store along with consumer and producer applications
  in ndnSIM, and reproduce the results of the base paper on the same tandem topology.
  \item Measure the retry delay of blind similarity caching when the cache contents persist across
  requests, and compare it with the analytical model.
  \item Design an Interest format that carries the minimum similarity accepted by the user and a
  bound on the number of retries, so that a cache answers only when its best match is acceptable.
  \item Design a forwarding strategy that chooses between answering locally and forwarding based on
  the measured upstream delay, and evaluate it under link congestion and under burst packet loss.
  \item Compare the proposed scheme with the base scheme in terms of mean and 95th percentile
  delivery delay, delivered similarity, upstream traffic and fairness across clients.
\end{enumerate}

The project will show whether simple changes to the request format and to the decision rule at
each cache can make similarity caching dependable in multi-hop networks, where a wrong answer from
a nearby cache costs the user a complete retry.

\subsection*{Methodology Overview}
The proposed system architecture consists of:
\begin{itemize}
  \item ns-3 with the ndnSIM module as the simulation platform
  \item Similarity-based Content Store replacing exact name lookup at each cache node
  \item Consumer applications that check the similarity and retry, and a producer acting as the
  origin server
  \item Custom forwarding strategy implementing the answer-or-forward decision
  \item Point-to-point links with configurable bandwidth, delay, queue size and burst loss
  \item Python and MATLAB scripts comparing the simulation results with the analytical model
\end{itemize}

\section*{Team Details}
\begin{itemize}
  \item \textbf{Somyak Priyadarshi Mohanta} (241CS257) - Implementation lead, Content Store and
  forwarding strategy, system architecture
  \item \textbf{Pranav Shaji} (241CSXXX) - Consumer and producer applications, simulation scenarios
  \item \textbf{Chris Tony} (241CSXXX) - Analytical model, performance analysis, documentation
\end{itemize}

\begin{thebibliography}{9}

\bibitem{ndn}
L. Zhang et al., ``Named Data Networking,'' \textit{ACM SIGCOMM Computer Communication Review},
vol. 44, no. 3, pp. 66-73, 2014.

\bibitem{neglia}
G. Neglia, M. Garetto, and E. Leonardi, ``Similarity Caching: Theory and Algorithms,'' in
\textit{IEEE/ACM Transactions on Networking}, vol. 30, no. 2, pp. 475-486, 2022.

\bibitem{vcache}
L. G. Schroeder et al., ``vCache: Verified Semantic Prompt Caching,'' arXiv:2502.03771, 2025.

\bibitem{nakamura}
R. Nakamura and N. Kamiyama, ``Analysis of Similarity Caching on General Cache Networks,'' in
\textit{IEEE Access}, vol. 12, pp. 163338-163348, 2024.

\bibitem{garetto}
M. Garetto, E. Leonardi, and G. Neglia, ``Content Placement in Networks of Similarity Caches,''
\textit{Computer Networks}, vol. 201, Art. no. 108570, 2021.

\bibitem{ndnsim}
S. Mastorakis, A. Afanasyev, and L. Zhang, ``On the Evolution of ndnSIM: an Open-Source Simulator
for NDN Experimentation,'' \textit{ACM SIGCOMM Computer Communication Review}, vol. 47, no. 3,
pp. 19-33, 2017.

\end{thebibliography}

\end{document}
